\section{Rational boxes and ambient crossing germs}\label{app:rational}

Our immediate objective is to continue ordinary local crossing along
rational trees, with enough information to retain a holomorphic germ on
the ambient manifold at every parameter pair. Throughout this appendix,
$M$ is a smooth projective complex manifold with
a holomorphic splitting
\[
T_M=E_1\oplus E_2
\]
into integrable subbundles. Write $\mathcal F_i$ for the corresponding
regular foliations. Frobenius submersions for the two foliations, taken
together, give holomorphic product charts $z=(z_1,z_2)$ in which $E_i$
is the $i$th coordinate distribution. Changes between these charts have,
locally, the form
\[
(z_1,z_2)\longmapsto\bigl(F(z_1),G(z_2)\bigr).
\]
Consequently the local crossing
\begin{equation}\label{rat:local-crossing}
\kappa(x,y)=z^{-1}\bigl(z_1(x),z_2(y)\bigr)
\end{equation}
is independent, as a germ along the diagonal, of the product chart.
Our eventual application, Theorem~\ref{thm:crossing}, requires continuation
of these \emph{ambient} germs, rather than just crossing values on curves.

\begin{definition}\label{rat:definition}
Let $A,B$ be locally path connected spaces. A continuous map
$R:A\times B\to M$ is a \emph{box} if, near every parameter pair, there
are a product neighborhood $A'\times B'$ and a product chart on its
image in which
\[
z\bigl(R(a,b)\bigr)=\bigl(r_1(a),r_2(b)\bigr).
\]
Thus first-parameter slices locally follow $\mathcal F_1$ plaques and
second-parameter slices locally follow $\mathcal F_2$ plaques. In
particular these slices are continuous into their leaves with the leaf
topology. A \emph{rational tree} is a connected nodal curve whose
irreducible components are copies of $\Pp^1$ and whose dual graph is a
finite tree.
\end{definition}

For continuous maps on products of rational trees, holomorphic on each
product of components, the box condition follows from the corresponding
separate tangencies. Indeed, in a product chart each wrong-coordinate
function is constant on every branch of its parameter slice; continuity
makes these constants agree at a node.

The local charts agree with those of Section~\ref{sec:prelim}; the
definition here also allows singular and merely topological parameter
spaces. We recall two elementary facts about holonomy in this setting. A
continuous plaque-wise path in an $\mathcal F_2$ leaf transports first
transverse coordinates by composing the first-coordinate changes in a
finite chain of product charts. This gives a biholomorphic germ between
first transversals. Reversing the path gives the inverse germ. The germ
is unchanged under a leafwise homotopy with fixed endpoints: subdivide
the compact homotopy square into pieces mapping into foliation charts,
and cancel the transverse changes around their boundaries. The same
statements hold with the two foliations interchanged.

\begin{lemma}[Germs supplied by a box]\label{rat:germs}
Let $T$ be a rational tree or an interval, let $R:T\times T\to M$ be a
box, and put $u(s)=R(s,s)$. There is a coherent family of holomorphic
map germs
\[
c_{s,t}:(M\times M,(u(s),u(t)))\longrightarrow(M,R(s,t))
\]
which agrees with ordinary crossing when $s=t$ and satisfies
\begin{equation}\label{rat:differential}
\begin{aligned}
\dd_1c_{s,t}(T_M)&\subset E_1,&\qquad \dd_1c_{s,t}(E_2)&=0,\\
\dd_2c_{s,t}(T_M)&\subset E_2,&\qquad \dd_2c_{s,t}(E_1)&=0.
\end{aligned}
\end{equation}
Here coherence means that near each parameter pair a single ambient
holomorphic map represents all the nearby germs. This family is unique
among coherent families agreeing with ordinary crossing on the diagonal.
\end{lemma}

\begin{proof}
Fix $s,t$ and a path $\gamma$ in $T$ from $s$ to $t$. Set
$p=u(s)$, $q=u(t)$ and $r=R(s,t)$, and choose product charts
$z,\zeta,w$ at $p,q,r$, respectively. The path $R(s,\gamma)$ runs
from $p$ to $r$ in an $\mathcal F_2$ leaf. Its holonomy gives a
first-coordinate germ $h_1:z_1\to w_1$. Similarly,
$R(\gamma^{-1},t)$ gives a second-coordinate germ
$h_2:\zeta_2\to w_2$. Define
\begin{equation}\label{rat:germ-formula}
c_{s,t}(x,y)=w^{-1}\bigl(h_1(z_1(x)),h_2(\zeta_2(y))\bigr).
\end{equation}
Chart changes cancel in this formula, and
\eqref{rat:differential} follows immediately. Both slices are continuous
into their leaves. Since $T$ is simply connected, two choices of
$\gamma$ give fixed-endpoint leafwise homotopies, so the formula is
independent of this choice.

For completeness, coherence concerns the entire formula, not merely its
value. Choose finite chains of product charts along
$R(s,\gamma)$ and $R(\gamma^{-1},t)$ realizing $h_1,h_2$. For $s'$ near
$s$ and $t'$ near $t$, continuity on these compact paths ensures that
the same chains apply to $R(s',\gamma)$ and $R(\gamma^{-1},t')$. Join
$s'$ to $s$ and $t$ to $t'$ inside small connected parameter
neighborhoods. All the resulting endpoint segments lie in the endpoint
product charts, so introduce no additional transverse-coordinate
change. The same holomorphic representative
\eqref{rat:germ-formula} therefore represents $c_{s',t'}$. At a node
one uses a small star neighborhood and connectors through the node;
no differentiability of the parameters is involved.

When $s=t$, choose the constant path to recover
\eqref{rat:local-crossing}. Finally, a coherent family is analytic
continuation along paths in the connected space $T\times T$ of its
germ at any chosen diagonal parameter. Uniqueness of continuation of
ambient holomorphic germs proves the asserted uniqueness.
\end{proof}

\begin{remark}\label{rat:coherence}
Coherence in Lemma~\ref{rat:germs} is local over the parameter space.
Distinct parameter pairs with the same ambient input need not yet give
the same germ. That identification is the purpose of the monodromy
argument in Theorem~\ref{thm:crossing}. Likewise, the value and
\eqref{rat:differential} alone do not specify an ambient germ.
\end{remark}

\begin{lemma}[Extension of a rational box]\label{rat:extension}
Put $S=\Pp^1\times\Pp^1$. A holomorphic box on a neighborhood of the
diagonal, or of the union
\[
D=(\Pp^1\times\{t_0\})\cup(\{s_0\}\times\Pp^1),
\]
extends uniquely to a holomorphic box on $S$.
\end{lemma}

\begin{proof}
In either case the divisor has class $(1,1)$ and is a hyperplane
section of the Segre embedding of $S$. Its complement $V$ is therefore
a smooth affine, hence Stein, surface. Shrink the given neighborhood
to a connected neighborhood $W$ of the divisor. Removing a complex
hypersurface does not disconnect $W$, so $W\setminus D$ is connected,
where here $D$ also denotes the diagonal in the first case. The set
$K=S\setminus W$ is compact in $V$, and $V\setminus K=W\setminus D$.
Meromorphic Hartogs extension on Stein surfaces, in the form of
\cite[Theorem~2.2]{MP09}, extends meromorphic functions from
$V\setminus K$ to $V$.

Embed $M$ in projective space and choose a coordinate hyperplane not
containing the image of the given map. Extend its projective coordinate
ratios and glue to their original values on $W$. They define a
meromorphic map $F:S\dashrightarrow M$: the equations of the embedded
$M$ hold on $W$ and hence identically. The two slice tangencies also
persist on the holomorphic locus by the identity theorem.

The meromorphic graph is a closed analytic subset of the projective
manifold $S\times M$, so it is algebraic by Chow's theorem
\cite[Proposition~13]{Serre56}. Thus $F$ is a rational map. Clearing
the poles of its finitely many projective coordinate ratios expresses
it by a finite-dimensional linear system on $S$. Remove the fixed
divisorial part of this system and resolve its base ideal by point
blowups. This is the classical surface elimination-of-indeterminacy argument
\cite{Beauville96}; the local intersection calculation gives termination
in the present setting.
Successively blow up the remaining basepoints: the local intersection number of two general
members of the resulting linear system drops by their positive
multiplicity product at each unresolved basepoint. The finite sum of
these intersection numbers therefore forces termination. This gives
$\mu:\widehat S\to S$ and a holomorphic map $H:\widehat S\to M$.
We show that every exceptional curve is contracted by $H$. At a generic
point of an exceptional component $C$ over $(s_0,t_0)$, choose
coordinates $(x,y)$ with $C=\{x=0\}$. For local product coordinates
$s,t$ downstairs, write
\[
s\circ\mu-s_0=x^a U(x,y),\qquad
t\circ\mu-t_0=x^b V(x,y),
\]
where $a,b>0$ and $U,V$ are units. The holomorphic vector fields
\[
W_s=-xU_y\partial_x+(aU+xU_x)\partial_y,
\qquad
W_t=-xV_y\partial_x+(bV+xV_x)\partial_y
\]
annihilate $\dd(s\circ\mu)$ and $\dd(t\circ\mu)$, respectively.
Off the exceptional locus, slice tangency gives
\[
\dd H(W_s)\in E_2,\qquad \dd H(W_t)\in E_1.
\]
These inclusions extend holomorphically to $C$. There both vector
fields are nonzero multiples of $\partial_y$, whence
\[
\dd H(T_C)\subset E_1\cap E_2=0.
\]
Thus $H$ is constant on each exceptional component. Fibers of $\mu$
are connected, so each whole fiber has a single image. Around this
image choose a target coordinate chart; properness gives a neighborhood
downstairs whose inverse image maps into the chart. Its coordinate
functions descend holomorphically, since
$\mu_*\mathcal O_{\widehat S}=\mathcal O_S$, equivalently by removal
of the isolated singularity downstairs. Hence $F$ is holomorphic.
Its slice tangencies give the box condition, and uniqueness follows
from the identity theorem.
\end{proof}

We have extended a box from one diagonal or one pair of axes. A rational
tree has a unique route between components; this lets us repeat the axes
construction without making a new compatibility choice.

\begin{proposition}[Rational boxes]\label{rat:boxes}
For every holomorphic map $u:T\to M$ from a rational tree there is a
unique continuous box
\[
R_u:T\times T\longrightarrow M,
\qquad R_u(s,s)=u(s),
\]
holomorphic on every product of irreducible components. If a rational
tree $T'$ contains $T$ as a subtree and $u':T'\to M$ extends $u$, then
$R_{u'}|_{T\times T}=R_u$. The ambient germ families of
Lemma~\ref{rat:germs} agree on this restriction as well.
\end{proposition}

\begin{proof}
First take $T=\Pp^1$. Applying local crossing to $u(s),u(t)$ gives
coherent holomorphic germs along the diagonal of $S$. Compactness
allows them to be glued on a neighborhood of the diagonal: take a
finite family of representatives, shrink their domains, and then
shrink around the diagonal so that overlapping representatives agree.
Lemma~\ref{rat:extension} gives a box on $S$. In a product chart any
box near $(s,s)$ has the form $(a(s),b(t))$; its diagonal specifies
both functions. This proves uniqueness near the diagonal and hence
on $S$.

We also need the corresponding construction from axes. Suppose
$\alpha:\Pp^1\to M$ is tangent to $E_1$,
$\beta:\Pp^1\to M$ is tangent to $E_2$, and
$\alpha(s_0)=\beta(t_0)$. Transport the germ of $\beta$ at $t_0$
along $\alpha$, using $\mathcal F_1$ holonomy. There is no path
ambiguity: a loop in the parameter sphere contracts, and its image
gives a leafwise homotopy. In product charts the resulting map near
$(s,t_0)$ has first coordinate supplied by $\alpha(s)$ and second
coordinate a holomorphic germ in $t$. It is jointly holomorphic and
the germs are coherent by the finite-chart argument above. Do the
symmetric construction along the other axis. Near their intersection
both maps are ordinary crossing, so, after shrinking, they glue on a
neighborhood of the axes. Lemma~\ref{rat:extension} extends the box
over $S$. The two prescribed axes determine its germ at $(s_0,t_0)$,
so the extension is unique. This includes constant curves.

For a general tree, first construct the box on the square of each
component. Proceed through ordered pairs of distinct components by
increasing distance in the dual tree. On $C_i\times C_j$, prescribe
the two axes at the nodes pointing toward each other along the unique
component route. These axes have already been defined at smaller
distances; they are tangent to $E_1,E_2$, respectively, and agree at
their corner by the previous edge identifications. Apply the axes
construction. Every common edge matches: for two adjacent components,
their distances to a fixed third component differ by one, and the
common edge is an inward axis for the later pair. The finitely many
component maps therefore glue continuously, and their separate
tangencies imply the box condition also at nodes.

The same induction proves uniqueness, first on component squares and
then from the inward axes. Restriction to a subtree is consequently
$R_u$. To compare its ambient germs, compute the two holonomies along
paths lying in $T$. The paths and box values are unchanged in $T'$,
so Equation~\eqref{rat:germ-formula} gives precisely the same germs.
\end{proof}
